% Part: computability
% Chapter: computability-theory
% Section: applications-fixed-point

\documentclass[../../../include/open-logic-section]{subfiles}

\begin{document}

% SEGMENT OLP-0250-S01

\olfileid{cmp}{thy}{apf}
\olsection{நிலைப்புள்ளித் தேற்றத்தைப் பயன்படுத்துதல்}

பகுதி கணிக்கத்தக்க சார்புகளை அவற்றின் சுட்டெண்களைக் கொண்டு
வரையறுக்க நிலைப்புள்ளித் தேற்றம் அடிப்படையில் அனுமதிக்கிறது.
எடுத்துக்காட்டாக, ஒவ்வொரு $y$ க்கும்
\[
\cfind{e}(y) = e + y.
\]
என அமையும் சுட்டெண் $e$ ஒன்றைக் கண்டுபிடிக்கலாம். இன்னோர்
எடுத்துக்காட்டாக, தன்னையே அச்சிடும் Java அல்லது C++ நிரலை
வடிவமைக்க நிலைப்புள்ளித் தேற்றத்தின் நிரூபணத்தைப் பயன்படுத்தலாம்.

ஒவ்வொரு $e$ க்கும் $\cfind{e}$ இன் வரையறைப்புலமாக $W_e$ ஐ
எடுத்துக்கொண்டால், தொடர் $W_0$, $W_1$, $W_2$,~\dots ஆனது
கணிக்கத்தக்கவாறு எண்ணிடத்தக்க கணங்களை எண்ணிடுகிறது என்பதை நினைவுகூர்க.
இக்கணங்களில் சில கணிக்கத்தக்கவை. மதிப்பு $x$ ஐ உள்ளீடாகப் பெற்று,
$W_x$ கணிக்கத்தக்கதாக இருக்க நேர்ந்தால் அதன் சிறப்பியல்புச்
சார்புக்கான ஒரு சுட்டெண்ணைத் திருப்பித்தரும் படிமுறை ஒன்று உள்ளதா
என்று கேட்கலாம். விடை ``இல்லை''; அத்தகைய படிமுறை இல்லை:

\begin{thm}
பின்வரும் பண்பைக் கொண்ட பகுதி கணிக்கத்தக்க சார்பு $f$ எதுவும் இல்லை:
$W_e$ கணிக்கத்தக்கதாக இருக்கும் போதெல்லாம், $f(e)$ வரையறுக்கப்பட்டு,
$\cfind{f(e)}$ அதன் சிறப்பியல்புச் சார்பாக இருக்கும்.
\end{thm}

\begin{proof}
$f$ ஏதாவது ஒரு கணிக்கத்தக்க சார்பு ஆகட்டும்; $W_e$ கணிக்கத்தக்கதாக
இருந்தும் $\cfind{f(e)}$ அதன் சிறப்பியல்புச் சார்பாக இல்லாதவாறு
ஓர்~$e$ ஐக் கட்டமைப்போம். நிலைப்புள்ளித் தேற்றத்தைப் பயன்படுத்தி,
\[
\cfind{e}(y) \simeq
\begin{cases}
0 & \text{$y=0$ மற்றும் $\cfind{f(e)}(0) \fdefined = 0$ எனில்} \\
\text{வரையறுக்கப்படவில்லை} & \text{இல்லையெனில்.}
\end{cases}
\]
என அமையும் சுட்டெண் $e$ ஒன்றைக் கண்டுபிடிக்கலாம். அதாவது,
பின்வருமாறு வரையறுக்கப்பட்ட சார்புக்கு நிலைப்புள்ளித் தேற்றத்தைப்
பயன்படுத்தி $e$ பெறப்படுகிறது:
\[
g(x,y) \simeq
\begin{cases}
0 & \text{$y=0$ மற்றும் $\cfind{f(x)}(0) \fdefined = 0$ எனில்} \\
\text{வரையறுக்கப்படவில்லை} & \text{இல்லையெனில்.}
\end{cases}
\]
$g$ பகுதி கணிக்கத்தக்கது என்பதை முறைசாராமல் பின்வருமாறு காணலாம்:
உள்ளீடுகள் $x$, $y$ இல், $y$ ஆனது~$0$ குச் சமமா என்று படிமுறை
முதலில் சரிபார்க்கிறது. சமம் எனில், படிமுறை $f(x)$ ஐக் கணித்து,
பிறகு அனைத்தளாவிய பொறியைக் கொண்டு $\cfind{f(x)}(0)$ ஐக் கணிக்கிறது.
இக்கடைசிக் கணிப்பு நின்று~$0$ ஐத் திருப்பித்தந்தால், படிமுறை~$0$ ஐத்
திருப்பித்தருகிறது; இல்லையெனில் படிமுறை நிற்பதில்லை.

இப்போது $\cfind{f(e)}(0)$ வரையறுக்கப்பட்டு $0$ குச் சமம் எனில்,
$y$ ஆனது $0$ குச் சமமாக இருக்கும்போது சரியாக $\cfind{e}(y)$
வரையறுக்கப்படுகிறது; எனவே $W_e = \{ 0 \}$. $\cfind{f(e)}(0)$
வரையறுக்கப்படவில்லை என்றாலோ, வரையறுக்கப்பட்டும் $0$ குச் சமமல்ல
என்றாலோ, $W_e = \emptyset$. இரு நேர்வுகளிலும் $\cfind{f(e)}$
ஆனது~$W_e$ இன் சிறப்பியல்புச் சார்பு அல்ல; ஏனெனில் உள்ளீடு $0$
இல் அது தவறான விடையைத் தருகிறது.
\end{proof}

\end{document}
